\documentclass[unicode,10pt,a4paper,oneside,numbers=endperiod,openany]{scrartcl}

\input{assignment.sty}
\usepackage{amsmath}
\begin{document}


\setassignment
\setduedate{09 March 2026, 18:00~(Europe/Zurich)}

\serieheader{High-Performance Computing Lab for CSE}{2026}
            {Student: Riccardo Polo}
            {Discussed with: FULL NAME}
            {Report for Project 1}{}
\newline

% you may comment this for the final submission, but make sure to comply with
% the policy
\assignmentpolicy

\newpage
\section{Euler warm-up [10 points]}
\textbf{Answer to Question 1.1} \\

\textbf{Answer to Question 1.2} \\

\textbf{Answer to Question 1.3} \\

\newpage
\section{Performance characteristics [50 points]}

\subsection{Peak performance [16 Points]}

Proceeding as in the example given in the project description, the hardware specifications can be obtained from the Euler webpage. \\
For Euler VII (phase 1), the system consists of $n_{\text{nodes}} = 292$ nodes, each equipped with $n_{\text{sockets}} = 2$ sockets and $n_{\text{cores}} = 64$ cores per socket, operating at a frequency $f = 2.6\,\text{GHz}$.\\
 For Euler VII (phase 2), the system consists of $n_{\text{nodes}} = 248$ nodes, each with $n_{\text{sockets}} = 2$ sockets and $n_{\text{cores}} = 64$ cores per socket, operating at a frequency $f = 2.45\,\text{GHz}$.\\

Using the uops.info database provided, we determine both the superscalability factor and the FMA factor for the corresponding AMD architectures. \\
For both Euler VII phases, which are based on AMD Zen2 and AMD Zen3 respectively, we obtain $n_{\text{super}} = 4$ and $n_{\text{FMA}} = 2$.\\

The FMA factor is determined by observing that all instructions marked as \textbf{FMA} in the table have a corresponding number of micro-operations equal to one. This indicates that the fused multiply-add is executed as a single micro-operation by dedicated hardware. Since an FMA instruction performs the operation $a \cdot b + c$, it consists of one floating-point multiplication and one floating-point addition. Therefore, each FMA instruction corresponds to two floating-point operations, and thus $n_{\text{FMA}} = 2$.

The superscalability is defined as the number of floating-point operations that the FPU of a core can execute in parallel within a single clock cycle. To determine this value, we analyze the instruction \texttt{VFMADD132PD (YMM, YMM, YMM)} reported in the uops.info table. The table specifies a throughput of $0.50$ cycles per instruction, one micro-operation, and execution on ports FP0/1. A throughput of $0.50$ cycles per instruction implies that

$$
\frac{1}{0.50} = 2
$$

FMA instructions can be issued per cycle. The port information FP0/1 confirms the presence of two independent floating-point execution pipelines capable of executing FMA instructions in parallel.

Since two FMA instructions can be executed per cycle and each FMA instruction corresponds to two floating-point operations, the total number of floating-point operations per cycle is

$$
2 \times 2 = 4.
$$

Hence, for both AMD Zen~2 and AMD Zen~3, the superscalability is

$$
4 \text{ floating-point operations per cycle per core}.
$$

To determine $n_{\text{SIMD}}$, we inspect the floating-point instructions reported in the uops.info table, where instructions such as \texttt{VFMADD132PD (YMM, YMM, YMM)} appear. The presence of YMM operands indicates the use of 256-bit AVX registers, since YMM registers are defined as 256-bit wide in the AVX/AVX2 instruction set. As both AMD Zen~2 and Zen~3 execute these YMM-based instructions, they necessarily support AVX2 with 256-bit SIMD registers. A double-precision floating-point number occupies 64 bits; therefore, the number of doubles that can be processed concurrently by one SIMD instruction is

$$
n_{\text{SIMD}} = \frac{256}{64} = 4.
$$

Hence, each SIMD instruction can operate on four double-precision values in parallel.

\begin{table}[h]
\centering
\begin{tabular}{lcc}
\hline
\textbf{Parameter} & \textbf{Euler VII (Phase 1)} & \textbf{Euler VII (Phase 2)} \\
\hline
$n_{\text{nodes}}$ & 292 & 248 \\
$n_{\text{sockets}}$ & 2 & 2 \\
$n_{\text{cores}}$ (per socket) & 64 & 64 \\
Clock frequency $f$ & $2.6\,\text{GHz}$ & $2.45\,\text{GHz}$ \\
$n_{\text{FMA}}$ & 2 & 2 \\
$n_{\text{super}}$ (FP ops/cycle/core) & 4 & 4 \\
$n_{\text{SIMD}}$ (doubles per SIMD instr.) & 4 & 4 \\
\hline
\end{tabular}
\caption{Summary of hardware parameters and architectural factors used for peak performance estimation.}
\end{table}
We can now compute all the peak performance parameters for both Euler VII phases :
\begin{align*}
&P_{\text{core,1}} =  n_{\text{super}} \times  n_{\text{FMA}} \times n_{\text{SIMD}} \times f = 4 \times 2 \times 4 \times 2.6\,\text{GHz} = 83.2\,\text{GFLOPS/core}. \\
&P_{\text{CPU,1}} = n_{\text{cores}} \times P_{\text{core,1}} = 64 \times 83.2\,\text{GFLOPS/core} = 5.3248\,\text{TFLOPS/CPU}. \\
&P_{\text{node,1}} = n_{\text{sockets}} \times P_{\text{CPU,1}} = 2 \times 5.3248\,\text{TFLOPS/node} = 10.6496\,\text{TFLOPS/node}. \\
&P_{\text{cluster,1}} = n_{\text{nodes}} \times P_{\text{node,1}} = 292 \times 10.6496\,\text{TFLOPS/node} = 3110.08\,\text{TFLOPS}.
\end{align*}


\begin{align*}
&P_{\text{core,2}} =  n_{\text{super}} \times  n_{\text{FMA}} \times n_{\text{SIMD}} \times f = 4 \times 2 \times 4 \times 2.45\,\text{GHz} = 78.4\,\text{GFLOPS/core}. \\
&P_{\text{CPU,2}} = n_{\text{cores}} \times P_{\text{core,2}} = 64 \times 78.4\,\text{GFLOPS/core} = 5.017  \,\text{TFLOPS/CPU}. \\
&P_{\text{node,2}} = n_{\text{sockets}} \times P_{\text{CPU,2}} = 2 \times 5.017\,\text{TFLOPS/node} = 10.034\,\text{TFLOPS/node}. \\
&P_{\text{cluster,2}} = n_{\text{nodes}} \times P_{\text{node,2}} = 248 \times 10.034\,\text{TFLOPS/node} = 2488.432\,\text{TFLOPS}.
\end{align*}

\subsection{Memory Hierarchies [8 Points]}
% template for the memory hierarchy table (adapt accordingly!)
% \begin{table}[hbt]
%   \centering
%   \begin{tabular}{| l | c |}
%     \hline
%     Main memory \hspace{1.5cm} &  32 GB \\
%     \hline
%     L4 cache & 128 MB \\
%     \hline
%     L3 cache & 6 MB \\
%     \hline
%     L2 cache & 256 KB \\
%     \hline
%     L1 cache & 32 KB \\
%     \hline
%   \end{tabular}
%   \caption{Memory hierarchy of an Euler login node with an Intel Xeon
%            E3-1284Lv4 quad-core CPU.
%            Each core has its own L1 and L2 cache, while L3 and L4 caches are
%            shared among all the four cores.}
%   \label{tab:perf_mem}
% \end{table}

% template for the memory topology figure (adapt accordingly!)
% \begin{figure}[h]
%   \centering
%   \includegraphics[width=0.5\textwidth]{images/XEON_E3-1284Lv4}
%   \caption{Topology of Euler login node with an Intel Xeon E3-1284Lv4 quad-core
%            CPU.}
%   \label{fig:hwloc-ls_XEON_E3-1284Lv4}
% \end{figure}

\newpage
\subsubsection{Cache and main memory size}
ripolo@eu-login-20:~$ lscpu
Architecture:                x86_64
  CPU op-mode(s):            32-bit, 64-bit
  Address sizes:             39 bits physical, 48 bits virtual
  Byte Order:                Little Endian
CPU(s):                      8
  On-line CPU(s) list:       0-7
Vendor ID:                   GenuineIntel
  Model name:                Intel(R) Xeon(R) CPU E3-1585L v5 @ 3.00GHz
    CPU family:              6
    Model:                   94
    Thread(s) per core:      2
    Core(s) per socket:      4
    Socket(s):               1
    Stepping:                3
    CPU max MHz:             3700.0000
    CPU min MHz:             800.0000
    BogoMIPS:                6000.00
    Flags:                   fpu vme de pse tsc msr pae mce cx8 apic sep mtrr pge mca cmov pat pse36 clflush dts acpi mm
                             x fxsr sse sse2 ss ht tm pbe syscall nx pdpe1gb rdtscp lm constant_tsc art arch_perfmon peb
                             s bts rep_good nopl xtopology nonstop_tsc cpuid aperfmperf pni pclmulqdq dtes64 monitor ds_
                             cpl vmx smx est tm2 ssse3 sdbg fma cx16 xtpr pdcm pcid sse4_1 sse4_2 x2apic movbe popcnt ts
                             c_deadline_timer aes xsave avx f16c rdrand lahf_lm abm 3dnowprefetch cpuid_fault epb invpci
                             d_single pti ssbd ibrs ibpb stibp tpr_shadow vnmi flexpriority ept vpid ept_ad fsgsbase tsc
                             _adjust bmi1 avx2 smep bmi2 erms invpcid mpx rdseed adx smap clflushopt intel_pt xsaveopt x
                             savec xgetbv1 xsaves dtherm ida arat pln pts hwp hwp_notify hwp_act_window hwp_epp md_clear
                              flush_l1d arch_capabilities ibpb_exit_to_user
Virtualization features:
  Virtualization:            VT-x
Caches (sum of all):
  L1d:                       128 KiB (4 instances)
  L1i:                       128 KiB (4 instances)
  L2:                        1 MiB (4 instances)
  L3:                        8 MiB (1 instance)


\subsection{Bandwidth: STREAM benchmark [10 Points]}


\subsection{Performance model: A simple roofline model [16 Points]}


\section{Auto-vectorization [10 points]}


\section{Matrix multiplication optimization [30 points]}


\end{document}
